% ============================================================
%  Evaluation & Conclusion – IncludeOS VirtioFS
%  Master's Thesis
% ============================================================

\section{Evaluation}
\label{sec:evaluation}

This section presents the experimental evaluation of the VirtioFS
implementation for IncludeOS.  Three research questions are investigated
through a set of controlled microbenchmarks and an end-to-end application
workload.  Section~\ref{sec:setup} describes the common infrastructure
shared by all experiments.  Sections~\ref{sec:rq1}–\ref{sec:rq3} then
present the setup, results, and analysis for each research question in
turn.

% -------------------------------------------------------------------
\subsection{Experimental Setup}
\label{sec:setup}

\subsubsection{Virtual-Machine Infrastructure}

All experiments are executed inside virtual machines managed by QEMU/KVM
on the same host workstation, ensuring that CPU microarchitecture and
memory topology are constant across every run.  A shared host directory
is exposed to the guest through VirtioFS, the primary transport under
evaluation; no virtual block device is involved.  The host directory
itself is mounted on a \texttt{tmpfs} ramdisk so that host-side storage
latency and caching effects are eliminated from the measurements.

Two guest operating environments are used:

\begin{description}
  \item[IncludeOS guest.]  The IncludeOS unikernel is launched directly
    by the \texttt{boot} utility (part of the IncludeOS toolchain), which
    invokes \texttt{qemu-system-x86\_64} with hardware-virtualisation
    (\texttt{-enable-kvm}) and 3\,GiB of memory.  No full operating
    system is present: the benchmark binary \emph{is} the unikernel image,
    linked with the IncludeOS runtime and its VirtioFS driver.  The
    VirtioFS device is configured directly through a JSON configuration
    file passed to \texttt{boot}.

  \item[Linux guest.]  An Ubuntu 24.04 server image runs inside QEMU with
    KVM enabled, 8\,GiB of memory, and four virtual CPUs (\texttt{-smp 4}).
    Host memory is allocated via \texttt{memfd} and shared with the guest
    (\texttt{memory-backend-memfd, share=on}).  The VirtioFS device is
    backed by the userspace \texttt{virtiofsd} daemon
    (\texttt{vhost-user-fs-pci}), launched on the host before QEMU starts.
    \texttt{virtiofsd} is configured with direct I/O
    (\texttt{-{}-allow-direct-io}) and caching disabled
    (\texttt{-{}-cache never}) to prevent daemon-side buffering from
    masking driver latency.
\end{description}

Two distinct notification mechanisms are exercised in the IncludeOS
experiments:

\begin{description}
  \item[Polling.]  After submitting a request to the VirtioFS virtqueue,
    the IncludeOS driver busy-polls the used-ring descriptor until the
    host returns a completion.  No interrupt is taken; the vCPU never
    halts.

  \item[IRQ (MSI-X).]  After submission, the vCPU executes a \texttt{HLT}
    instruction and sleeps until a MSI-X interrupt from the VirtioFS
    device signals completion.  This frees the vCPU for other work but
    introduces interrupt-delivery latency on every I/O.
\end{description}

\subsubsection{Reproducible Builds with Nix}

All benchmark programs are compiled with the Nix package manager to
guarantee bit-for-bit reproducibility across machines.  A top-level
\texttt{shell.nix} provides the build environment and runtime
dependencies—including \texttt{virtiofsd}, QEMU, GRUB, and the IncludeOS
VM-runner utility—through a single \texttt{nix-shell} invocation.

IncludeOS unikernel binaries are built with
\texttt{includeos.stdenv.mkDerivation}; the resulting
\texttt{virtiofs\_bench.elf.bin} ELF image is passed directly to the
\texttt{boot} launcher.  Linux binaries are compiled with the
\texttt{pkgsMusl} static C library toolchain to minimise runtime
dependencies and prevent dynamic linker behaviour from confounding
results.  The C++20 standard is required by all benchmarks for
\texttt{std::format} and \texttt{std::chrono::high\_resolution\_clock}.

The orchestration script \texttt{run\_bench.sh} automates the full
workflow: it mounts the shared \texttt{tmpfs}, copies benchmark material
into it, invokes the appropriate Nix build, starts the virtual machine,
waits for it to finish, copies CSV result files back to the host, and
cleans up.

% -------------------------------------------------------------------
\subsubsection{POSIX I/O Benchmark (Experiments 1 \& 2)}
\label{sec:posix_setup}

The synchronous-I/O benchmark exercises VirtioFS through the standard
POSIX file API (\texttt{open}/\texttt{read}/\texttt{write}/\texttt{lseek}/\texttt{close}).
In IncludeOS this translates directly into driver calls with no kernel
syscall boundary; in Linux the calls traverse the VFS and FUSE layers
before reaching the \texttt{virtiofsd}-backed driver.

\paragraph{Test file.}  A 204.8\,MiB binary file
(\texttt{syncio\_testing\_file.bin}) pre-placed on the shared
\texttt{tmpfs} serves as the read source.  Write benchmarks create a
fresh output file on each run and remove it afterwards to prevent caching
warm-up from propagating across runs.  Files are opened with
\texttt{O\_DIRECT} to bypass the page cache, ensuring that all data
transfers reach the VirtioFS transport layer.

\paragraph{Chunk sweep.}  The benchmark iterates over chunk sizes from
8\,KiB to 248\,KiB in steps of 8\,KiB (31 sizes).  For each chunk size
30 independent runs are executed.  Four access patterns are measured:
\emph{sequential read}, \emph{sequential write}, \emph{random read}, and
\emph{random write}.  The random benchmarks permute chunk indices with a
fixed seed (Mersenne Twister, seed 42) so that the same access order is
reproduced in every run.

\paragraph{Timing.}  Each run is timed with
\texttt{std::chrono::high\_resolution\_clock}, wrapping the innermost
processing loop.  File open and close operations are excluded.  The result
is recorded as elapsed time in milliseconds and as throughput in
mebibytes per second (MiB/s), computed as $\text{filesize} /
\text{elapsed}$.

\paragraph{Result files.}  Four CSV files are written to the shared
VirtioFS directory (\texttt{seq\_read\_bench.csv},
\texttt{seq\_write\_bench.csv}, \texttt{rng\_read\_bench.csv},
\texttt{rng\_write\_bench.csv}).  Each row contains
\texttt{chunk\_size}, \texttt{run\_number}, \texttt{time\_ms}, and
\texttt{throughput\_mibps}.

\paragraph{Verification.} With the \texttt{VERIFY} compile-time flag
enabled, the benchmark writes the final read buffer and write buffer back
to the shared directory as \texttt{*\_copy.bin} files.  These can be
compared with the original \texttt{syncio\_testing\_file.bin} via SHA-256
to confirm data integrity across all driver paths.

% -------------------------------------------------------------------
\subsubsection{io\_uring Benchmark (Experiment 3)}
\label{sec:iouring_setup}

The io\_uring benchmark uses the same access patterns and file sizes as
the POSIX benchmark but replaces the synchronous system calls with the
\texttt{liburing} asynchronous interface configured with the
\texttt{IORING\_SETUP\_SQPOLL} flag.  SQPOLL spawns a kernel thread that
polls the submission queue without requiring a \texttt{io\_uring\_enter}
system call from the application, theoretically reducing syscall overhead
at high submission rates.  The submission queue depth is set to 1; a
single outstanding I/O is submitted and waited upon before the next is
issued, matching the serialised access pattern of the POSIX benchmark and
making throughput comparisons directly meaningful.  This benchmark runs
only on Linux; IncludeOS does not support io\_uring.

% -------------------------------------------------------------------
\subsubsection{c63 Video Decoder Benchmark}
\label{sec:c63dec_setup}

The c63 codec is a simplified MPEG-style video codec developed for
educational purposes.  The decoder \texttt{c63dec} reads a compressed
\texttt{.c63} file from the VirtioFS-mounted directory and writes the
decoded raw YUV frames back to the same directory, exercising both the
read and write paths of the VirtioFS driver under a realistic,
compute-bound workload.

The input file is a 51-frame WWE broadcast stream (\texttt{wwe.c63}) of
limited resolution; the file is pre-copied to the shared \texttt{tmpfs}.
The entire decode-and-write sequence is timed with
\texttt{std::chrono::high\_resolution\_clock} and repeated 30 times per
system configuration.  Run~0 is treated as a cold-start observation and
excluded from statistical analysis; runs~1–29 are used to compute means
and confidence intervals.

Three system configurations are compared: Linux, IncludeOS Polling, and
IncludeOS IRQ.  The IncludeOS build uses the same Nix derivation as the
POSIX benchmark but without the \texttt{-DCMAKE\_CXX\_STANDARD} flag
override (the C++20 standard is implied by IncludeOS's default
toolchain).  The Linux build uses the \texttt{pkgsMusl} static toolchain.

% -------------------------------------------------------------------
\subsection{RQ1 – VirtioFS Throughput: IncludeOS vs.\ Linux}
\label{sec:rq1}

\textbf{RQ1:} \emph{Can a VirtioFS implementation in IncludeOS have a
higher I/O throughput than Linux?}

\subsubsection{Results}

Figure~\ref{fig:rq1_combined} shows the mean throughput (with 95\%
confidence intervals as shaded bands) across all 31 chunk sizes for the
four access patterns, comparing Linux POSIX, IncludeOS Polling, and
IncludeOS IRQ.  Figure~\ref{fig:rq1_peak} summarises peak throughput as
a bar chart.

\begin{figure}[htbp]
  \centering
  \includegraphics[width=\linewidth]{plots/rq1_combined.pdf}
  \caption{RQ1 – Mean VirtioFS throughput (MiB/s) vs.\ chunk size for
    all four access patterns.  Shaded bands represent 95\% confidence
    intervals over 30 runs.}
  \label{fig:rq1_combined}
\end{figure}

\begin{figure}[htbp]
  \centering
  \includegraphics[width=\linewidth]{plots/rq1_peak_bar.pdf}
  \caption{RQ1 – Peak throughput comparison for sequential (left) and
    random (right) I/O across the three system configurations.}
  \label{fig:rq1_peak}
\end{figure}

\paragraph{IncludeOS Polling vs.\ Linux POSIX.}
IncludeOS Polling achieves the highest throughput across every access
pattern and chunk size.  At peak (248\,KiB chunks for reads, around
192--248\,KiB for writes), it sustains
\textbf{4{,}786\,MiB/s} for sequential reads and
\textbf{2{,}076\,MiB/s} for sequential writes—representing speedups of
$\times$\textbf{1.74} and $\times$\textbf{1.84} over the Linux POSIX
baseline respectively.  Random access exhibits nearly identical relative
gains: $\times$1.78 for random reads and $\times$1.83 for random writes.

The advantage is even more pronounced at small chunk sizes, where system
call and interrupt overhead dominate.  At 8\,KiB chunks, IncludeOS
Polling delivers \textbf{804\,MiB/s} for sequential reads against only
\textbf{266\,MiB/s} for Linux POSIX—a $\times$\textbf{3.02} speedup.
This highlights how the elimination of the kernel-mode/user-mode
transition, the VFS dispatch, and the FUSE/\texttt{virtiofsd} round-trip
dramatically reduces per-request latency when requests are frequent and
small.

\paragraph{IncludeOS IRQ vs.\ Linux POSIX.}
IncludeOS IRQ outperforms Linux POSIX across the entire chunk range, but
by a narrower margin than polling.  Peak sequential read throughput
reaches \textbf{3{,}454\,MiB/s} ($\times$1.26 over Linux POSIX), and
peak sequential write reaches \textbf{1{,}696\,MiB/s} ($\times$1.50).
At 8\,KiB chunks the IRQ mode achieves $\times$1.24 over Linux for
sequential reads—substantially better than Linux but well below polling's
$\times$3.02.

The performance gap between IRQ and Polling within IncludeOS is itself
instructive.  Polling outperforms IRQ by $\times$1.39 at peak sequential
read and $\times$1.22 at peak sequential write.  This gap is attributable
to the round-trip cost of MSI-X delivery: after every \texttt{HLT}, the
vCPU must be woken from its halted state by the hypervisor before
execution can resume, a path that involves an exit from the guest,
interrupt injection, and a re-entry—all at a cost invisible to the
polling driver, which simply spins in user space.

\paragraph{Throughput scaling with chunk size.}
All three systems exhibit a consistent pattern: throughput rises steeply
from 8\,KiB and plateaus beyond approximately 128--192\,KiB, after which
there is little additional gain from larger transfers.  This plateau
reflects the saturation of the underlying \texttt{tmpfs} memory bandwidth
rather than a driver-level bottleneck.  The plateau for IncludeOS
Polling, however, sits roughly twice as high as that for Linux POSIX,
confirming that the unikernel's reduced software overhead shifts the
bottleneck earlier along the chunk-size axis.

\paragraph{Result variance.}
Both IncludeOS configurations exhibit markedly lower coefficient of
variation (CV) than Linux POSIX.  For sequential reads, IncludeOS
Polling has CV = 1.40\% and IncludeOS IRQ has CV = 1.70\%, compared with
CV = 5.31\% for Linux POSIX.  The reduced variance reflects the
deterministic execution model of the unikernel: there is no scheduler
preemption, no page-fault handling, and no background kernel activity
competing for CPU cycles.

\subsubsection{Analysis}

The throughput advantage of IncludeOS over Linux can be attributed to
three structural differences:

\begin{enumerate}
  \item \textbf{No syscall overhead.}  In IncludeOS, a call to
    \texttt{read()} or \texttt{write()} dispatches directly to the
    VirtioFS driver without a privilege-level switch.  In Linux, every
    POSIX call traverses the \texttt{syscall} entry stub, VFS
    \texttt{vfs\_read}/\texttt{vfs\_write}, and then the FUSE file-system
    layer, which forwards the request to \texttt{virtiofsd} in user space
    via a shared ring buffer—itself requiring a further context switch.
    The round-trip involves at minimum two mode switches and a user-space
    daemon wakeup per I/O operation.

  \item \textbf{No FUSE protocol overhead.}  Linux VirtioFS is
    implemented on top of the kernel FUSE framework.  Every data transfer
    is wrapped in a FUSE request header, routed through the FUSE device,
    and unpacked by \texttt{virtiofsd}.  IncludeOS communicates directly
    with the VirtioFS virtqueue without any FUSE framing.

  \item \textbf{Absence of OS interference.}  As a unikernel, IncludeOS
    runs a single cooperative task.  There are no timer interrupts, no
    background kernel threads, no memory-pressure events, and no
    scheduler preemption during the benchmark, all of which contribute to
    the lower variance and higher average throughput observed.
\end{enumerate}

\textbf{Answer to RQ1:} Yes.  The IncludeOS VirtioFS implementation
achieves substantially higher I/O throughput than Linux across all
measured access patterns, chunk sizes, and notification modes.  In
polling mode the peak speedup reaches $\times$1.74--1.84 at large chunk
sizes and $\times$2.34--3.02 at 8\,KiB—demonstrating that the
elimination of kernel, VFS, and FUSE overheads is highly beneficial for
VirtioFS-backed I/O.

% -------------------------------------------------------------------
\subsection{RQ2 – Tweaking Linux: io\_uring SQPOLL}
\label{sec:rq2}

\textbf{RQ2:} \emph{Why is the I/O throughput in VirtioFS Linux slower
than IncludeOS?  Can we tweak Linux to achieve equivalent performance?}

\subsubsection{Results}

Figure~\ref{fig:rq2_combined} compares Linux POSIX and Linux io\_uring
SQPOLL across the four access patterns.  Figure~\ref{fig:rq2_speedup}
shows the per-chunk-size speedup ratio of io\_uring relative to POSIX.

\begin{figure}[htbp]
  \centering
  \includegraphics[width=\linewidth]{plots/rq2_combined.pdf}
  \caption{RQ2 – Mean VirtioFS throughput: Linux POSIX vs.\ Linux
    io\_uring SQPOLL.  Shaded bands are 95\% confidence intervals.}
  \label{fig:rq2_combined}
\end{figure}

\begin{figure}[htbp]
  \centering
  \includegraphics[width=\linewidth]{plots/rq2_speedup.pdf}
  \caption{RQ2 – io\_uring SQPOLL throughput as a ratio of Linux POSIX
    throughput (values above 1.0 indicate an improvement, below 1.0 a
    regression).  The dashed grey line marks parity.}
  \label{fig:rq2_speedup}
\end{figure}

Contrary to what one might expect from a zero-copy, polling-based I/O
interface, io\_uring SQPOLL is \emph{slower} than Linux POSIX across the
entire chunk-size range and all four access patterns.  Peak sequential
read throughput with io\_uring reaches only \textbf{2{,}280\,MiB/s}
versus \textbf{2{,}743\,MiB/s} for POSIX (a regression of
$\times$0.83).  For sequential writes the regression is $\times$0.88.
At the smallest chunk size of 8\,KiB, the penalty is even larger:
io\_uring achieves only \textbf{166\,MiB/s} for sequential reads,
against \textbf{266\,MiB/s} for POSIX ($\times$0.62), indicating that
the overhead is particularly acute at high request rates.

Figure~\ref{fig:all_systems} places all four configurations side by side
to contextualise these findings relative to the IncludeOS baselines.

\begin{figure}[htbp]
  \centering
  \includegraphics[width=\linewidth]{plots/all_systems_combined.pdf}
  \caption{All system configurations – VirtioFS throughput overview for
    all four access patterns.  io\_uring falls below even the Linux POSIX
    baseline.}
  \label{fig:all_systems}
\end{figure}

\subsubsection{Analysis}

The degradation of io\_uring SQPOLL relative to POSIX on a VirtioFS
filesystem has several interacting causes:

\paragraph{VirtioFS is not a native io\_uring target.}  The Linux
VirtioFS driver sits inside the kernel's FUSE layer, which ultimately
forwards requests to \texttt{virtiofsd} in user space.  io\_uring
submits I/O operations to the kernel, but for FUSE filesystems the kernel
must still block waiting for the daemon to process the request.
Consequently, even though the application avoids the \texttt{read}/\texttt{write}
syscall entry through SQPOLL, the kernel I/O path itself remains
synchronous with respect to the daemon: io\_uring does not provide a
path to bypass the FUSE round-trip.

\paragraph{SQPOLL thread contention.}  The \texttt{IORING\_SETUP\_SQPOLL}
flag spawns a dedicated kernel thread that continuously polls the
submission queue.  This thread consumes a hardware thread at all times,
including when there is no outstanding I/O.  With \texttt{sq\_thread\_idle}
set to 10\,000\,ms the SQPOLL thread never sleeps during the benchmark.
On the Linux guest, which has four vCPUs, this persistent spinning
competes with the \texttt{virtiofsd} processing thread and the
application thread itself—increasing cache pressure and reducing the
effective CPU available to forward I/O.

\paragraph{Queue depth of 1.}  The benchmark deliberately maintains a
single outstanding I/O at a time to match the serialised POSIX access
pattern.  SQPOLL provides its greatest benefit at high queue depths where
many I/O operations can be batched into the submission ring with a single
kernel notification.  At depth 1 the overhead of the ring management
(submission queue entry formatting, memory barriers, and completion-queue
polling) is incurred for every single operation without amortisation.

\paragraph{Structural gap remains.}  Even if SQPOLL overhead were
eliminated, the FUSE protocol and \texttt{virtiofsd} round-trip would
remain.  The root causes of Linux's throughput gap relative to IncludeOS
are architectural: the kernel VFS layer, the FUSE framing, the
\texttt{virtiofsd} user-space daemon, and the privilege transitions that
separate them are all absent in a unikernel.  These cannot be removed
through a user-space API change alone; they would require kernel
modifications (e.g., a native kernel-space VirtioFS driver operating
outside FUSE) or a unikernel-like architecture.

\textbf{Answer to RQ2:}  The throughput gap between Linux and IncludeOS
is caused by the multiple software layers interposed between the
application and the VirtioFS virtqueue in a general-purpose OS: system
calls, the VFS, FUSE framing, and the \texttt{virtiofsd} user-space
daemon.  Replacing synchronous POSIX calls with io\_uring SQPOLL does
not close this gap; on the contrary, at queue depth~1 with a FUSE-backed
filesystem, SQPOLL introduces a throughput regression of 9--38\%
compared with ordinary POSIX I/O.  Achieving parity with IncludeOS on
Linux would require removing the FUSE indirection, which is a kernel-level
architectural change beyond the scope of a user-space API tweak.

% -------------------------------------------------------------------
\subsection{RQ3 – Impact on Video Decoding}
\label{sec:rq3}

\textbf{RQ3:} \emph{What impact will enhanced I/O throughput have on
common tasks such as video decoding?}

\subsubsection{Results}

Figure~\ref{fig:rq3_box} shows a box plot of decode times (runs 1--29)
for the three system configurations.  Figure~\ref{fig:rq3_bar} presents
means with 95\% confidence intervals and speedup annotations relative to
Linux.

\begin{figure}[htbp]
  \centering
  \includegraphics[width=0.8\linewidth]{plots/rq3_c63dec_boxplot.pdf}
  \caption{RQ3 – Distribution of c63 decode times over 29 warm runs.
    Boxes span the interquartile range; whiskers extend to the most
    extreme non-outlier observation.}
  \label{fig:rq3_box}
\end{figure}

\begin{figure}[htbp]
  \centering
  \includegraphics[width=0.7\linewidth]{plots/rq3_c63dec_bar.pdf}
  \caption{RQ3 – Mean c63 decode time ± 95\% CI.  Speedup labels above
    the IncludeOS bars are relative to the Linux mean.}
  \label{fig:rq3_bar}
\end{figure}

Table~\ref{tab:c63dec_stats} summarises the per-configuration statistics.

\begin{table}[htbp]
  \centering
  \caption{c63dec benchmark statistics (runs 1--29; run 0 excluded as
    cold start).  Speedup is relative to Linux POSIX.}
  \label{tab:c63dec_stats}
  \begin{tabular}{lrrrrc}
    \hline
    \textbf{Configuration} & \textbf{Mean (s)} & \textbf{Std (s)} &
      \textbf{Min (s)} & \textbf{Max (s)} & \textbf{Speedup} \\
    \hline
    Linux POSIX          & 4.861 & 0.032 & 4.820 & 4.957 & $\times$1.00 \\
    IncludeOS Polling    & 3.759 & 0.015 & 3.731 & 3.790 & $\times$\textbf{1.29} \\
    IncludeOS IRQ (MSI-X)& 4.096 & 0.030 & 4.020 & 4.144 & $\times$1.19 \\
    \hline
  \end{tabular}
\end{table}

\paragraph{IncludeOS Polling.}
IncludeOS Polling achieves a mean decode time of \textbf{3.759\,s},
compared with \textbf{4.861\,s} for Linux POSIX—a speedup of
$\times$\textbf{1.29}, or a \textbf{22.7\% reduction} in wall-clock
decode time.  This result is consistent with the throughput advantage
observed in the POSIX microbenchmarks: the c63dec workload performs
many small reads (parsing the compressed bitstream) and writes (emitting
decoded YUV frames), both of which benefit disproportionately from the
elimination of per-request POSIX and FUSE overhead.

IncludeOS Polling also exhibits the lowest variance of any configuration:
a standard deviation of just \textbf{0.015\,s} (CV~=~0.4\%) versus
\textbf{0.032\,s} (CV~=~0.7\%) for Linux.  The tighter distribution
makes performance more predictable—a property valuable in real-time or
latency-sensitive media pipelines.

\paragraph{IncludeOS IRQ.}
IncludeOS IRQ reduces decode time to \textbf{4.096\,s}, a
$\times$\textbf{1.19} speedup (18.6\% reduction) over Linux.  The 8.5\%
gap between IRQ and Polling within IncludeOS reflects the interrupt-delivery
overhead noted in Section~\ref{sec:rq1}: for a codec that issues many
small I/O requests per decoded frame, the accumulated cost of MSI-X wakeup
cycles is significant.

\paragraph{Cold-start behaviour.}
Run~0 reveals a striking asymmetry between Linux and IncludeOS.  Linux
requires \textbf{6.092\,s} on the cold run—a \textbf{25.3\%} penalty
over its warm mean—caused by page-fault handling, dynamic loader
initialisation, and I/O scheduling setup.  IncludeOS Polling's cold run
takes \textbf{3.772\,s}, only \textbf{0.4\%} above its warm mean;
IncludeOS IRQ shows a similarly negligible \textbf{0.2\%} cold overhead.
This behaviour reflects the unikernel's single-process, pre-linked
execution model: there is no dynamic linker, no kernel process
initialisation, and no page-table population at startup.

\paragraph{Decode speed in context.}
The 29\% wall-clock improvement represents a meaningful gain in a
realistic media-processing scenario.  A pipeline decoding video at a
fixed frame rate would be able to sustain a higher rate or leave more
CPU headroom for subsequent processing stages (encoding, filtering,
network transmission) when running on IncludeOS.  Conversely, the
determinism improvement—a threefold reduction in standard deviation—
reduces worst-case latency, which matters for applications with
real-time playback constraints.

\subsubsection{Analysis}

The magnitude of the end-to-end speedup (22.7\%) is smaller than the
raw peak-throughput speedup seen in the microbenchmarks ($\times$1.74 for
sequential reads).  This is expected: the c63dec workload is not
I/O-bound.  A substantial fraction of execution time is spent in the CPU
decode pipeline (VLC decoding, IDCT, motion compensation), which is
identical across all configurations.  The measured speedup therefore
represents the I/O contribution to the total workload, scaled by the
Amdahl factor of the non-I/O fraction.

\textbf{Answer to RQ3:}  The higher I/O throughput of the IncludeOS
VirtioFS implementation translates directly into faster end-to-end
execution of a realistic video decoding task.  IncludeOS Polling
completes decoding \textbf{22.7\% faster} than Linux POSIX, with a
\textbf{threefold reduction in timing variance}.  IncludeOS IRQ delivers
a more modest but still meaningful 18.6\% improvement.  The unikernel
also eliminates the significant cold-start penalty observed on Linux,
making every execution faster and more predictable.

% -------------------------------------------------------------------
\subsection{Threats to Validity}
\label{sec:threats}

\paragraph{Internal validity.}
The \texttt{tmpfs} shared directory eliminates disk I/O from the
measurements but also places an upper bound on achievable throughput set
by the host's DRAM bandwidth.  Both Linux and IncludeOS are subject to
the same bound; the relative comparison remains valid, but absolute
numbers cannot be directly extrapolated to persistent storage scenarios.
The benchmark disables all caching in \texttt{virtiofsd} and uses
\texttt{O\_DIRECT} to prevent page-cache effects from masking driver
latency; however, it cannot control KVM's internal shadow-page-table
behaviour or QEMU's virtqueue polling interval.

\paragraph{External validity.}
The benchmark uses a single test file size (204.8\,MiB) and a fixed
random seed.  Workloads with smaller files, varying access patterns, or
concurrency (multiple I/O threads) may exhibit different relative
performance.  The c63dec workload is a low-resolution educational codec;
production video workloads with larger frames, higher frame rates, or
compressed formats (H.264, AV1) may be more or less I/O-bound.

\paragraph{Reproducibility.}
All experiments use Nix-pinned dependency graphs; the build and runtime
environments are fully reproducible.  The orchestration script
\texttt{run\_bench.sh} is deterministic and requires no manual
intervention between runs.  Nevertheless, results may differ on hardware
with different DRAM bandwidth, CPU microarchitecture, or hypervisor
implementation.

% ============================================================
\section{Conclusion}
\label{sec:conclusion}

This thesis has presented the design, implementation, and evaluation of a
VirtioFS driver for the IncludeOS unikernel.  The evaluation addressed
three research questions through a suite of controlled microbenchmarks
and an end-to-end video decoding workload.

Regarding \textbf{RQ1}, the results confirm that IncludeOS can achieve
significantly higher VirtioFS I/O throughput than a standard Linux guest.
In polling mode, peak throughput exceeds Linux POSIX by $\times$1.74--1.84
across sequential and random patterns, and by up to $\times$3.02 at
small chunk sizes where per-request overhead dominates.  Even the
interrupt-driven MSI-X mode outperforms Linux by $\times$1.26--1.53.
The advantage stems from the structural simplicity of the unikernel: the
absence of a kernel boundary, VFS dispatch, FUSE framing, and a
user-space daemon reduces every I/O to a near-direct transfer to the
VirtioFS virtqueue.

Regarding \textbf{RQ2}, the investigation into Linux-side optimisations
via io\_uring SQPOLL produced a counterintuitive result: SQPOLL
\emph{degrades} VirtioFS throughput by 9--38\% relative to ordinary POSIX
I/O.  The regression arises because the FUSE-backed VirtioFS kernel driver
cannot benefit from io\_uring's submission-batching model—each I/O still
incurs a synchronous round-trip to \texttt{virtiofsd}—while the SQPOLL
kernel thread introduces additional CPU contention.  Closing the
performance gap with IncludeOS on Linux would require a kernel-level
redesign of the VirtioFS stack, such as moving the driver out of the FUSE
layer and into a dedicated in-kernel implementation.

Regarding \textbf{RQ3}, the end-to-end c63 video decoder benchmark shows
that the I/O throughput advantage of IncludeOS translates into meaningful
application-level gains.  IncludeOS Polling completes decoding 22.7\%
faster than Linux with three times lower timing variance.  Notably,
IncludeOS completely eliminates the 25.3\% cold-start penalty observed on
Linux, making unikernel execution faster from the very first run—a
property particularly valuable in short-lived, function-as-a-service
deployments where cold-start latency directly affects user experience.

Taken together, these findings validate the motivation for implementing
VirtioFS in a unikernel: the removal of general-purpose OS abstractions
that are unnecessary in a single-application environment yields
substantial, measurable I/O performance improvements that propagate to
realistic application workloads.  Future work could explore higher
io\_uring queue depths to amortise SQPOLL overhead, multi-threaded
virtqueue access within IncludeOS, and persistent storage backends beyond
\texttt{tmpfs} to evaluate the driver under disk-latency constraints.
